%==============================================================
%  Chapter 2 -- Datasets
%  Source: datasets document (co-author). Rewritten as thesis prose;
%  the product-embedding construction is deferred to the embedding
%  chapter, where those vectors are actually consumed.
%==============================================================
\chapter{Datasets}
\label{ch:datasets}

Every method in this thesis is built on data collected from real user activity
on the IndiaMART platform. This chapter describes that data before any model is
introduced, so that later chapters can refer back to a single, consistent
account of what was observed and how it was prepared. Three datasets are used.
The first is a log of raw browsing interactions; the second is a transition
dataset derived from those logs, which is the primary training input for the
sequential models; and the third is a catalogue of product descriptions that
supplies content information about individual items. The first two capture how
users behave, while the third describes what the products are.

The interaction log and the transition dataset are complementary views of the
same activity. The log keeps every interaction event as it happened, whereas
the transition dataset re-expresses that activity as ordered product-to-product
movements, which is the form the recommendation models need. The remainder of
the chapter takes each dataset in turn and then describes the preprocessing
applied before training.

%--------------------------------------------------------------
\section{Click Stream Log (CSL) Dataset}
\label{sec:csl}

The Click Stream Log records the interactions a user generates while browsing
IndiaMART. These are \emph{implicit} signals: they come from ordinary browsing
behaviour rather than from ratings or any other feedback a user is asked to
provide. Implicit feedback of this kind is the norm in large industrial
recommender systems, where explicit ratings are available for only a small
fraction of users and cannot be relied upon at scale.

The log contains several kinds of events corresponding to different stages of
engagement --- among them product impressions, product clicks, and conversion
events --- so that together they trace a user's path from being shown an item to
acting on it. Because the data is drawn from genuine browsing sessions, it
carries the characteristics one expects of a production recommendation setting:
interactions are sparse relative to the size of the catalogue, user interests
shift over time, and the product inventory is large and heterogeneous.

Each record is tagged with a user identifier, the browsing platform, a
timestamp, product information, and other contextual attributes from the
session. Because events carry timestamps, they can be ordered in time, and it is
this chronological ordering that makes it possible to reconstruct the sequential
browsing behaviour from which the transition dataset is built.

%--------------------------------------------------------------
\section{Transition Dataset}
\label{sec:transition-data}

The click stream log faithfully records individual interactions, but it does not
by itself express the sequential relationship between products a user views one
after another. Since the models in this work are concerned precisely with those
product-to-product movements, the log is converted during preprocessing into a
separate \emph{transition dataset}. This dataset is stored on its own and is the
main source used for training.

Each row of the transition dataset represents one observed transition from a
source product to a destination product. Alongside the two products, each row
keeps contextual information --- the browsing platform, the category identifiers
of both products, a timestamp, and metadata describing the kind of transition.
Representing the data this way lets a model learn patterns in how users navigate
between products, rather than treating each interaction as an isolated event.

The timestamp on each transition is retained deliberately, because it is later
used to split the data into training and test sets chronologically. Keeping the
temporal order intact prevents information from later interactions leaking into
the training of a model that is meant to predict them, and it lets the models be
evaluated under conditions closer to real deployment, where only past behaviour
is available when a recommendation is made.

\subsection{Schema}
\label{sec:transition-schema}
Table~\ref{tab:transition-schema} lists the attributes of the transition
dataset. Between them, the product identifiers record the sequential movement
from one item to the next, the category identifiers carry the semantic grouping
of each product, and the platform field allows behaviour to be compared across
browsing modes. Taken together these fields give a compact but informative
description of user navigation suitable for training recommendation models over
a very large product catalogue.

\begin{table}[ht]
  \centering
  \caption{Schema of the transition dataset.}
  \label{tab:transition-schema}
  \begin{tabular}{ll}
    \toprule
    Attribute & Description \\
    \midrule
    \texttt{User\_ID}        & Identifier of the user who generated the transition. \\
    \texttt{MODID}           & Browsing platform / mode of access (e.g.\ website or mobile). \\
    \texttt{Timestamp}       & Time of the transition; used for chronological splitting. \\
    \texttt{ProdID1}         & Source product of the transition. \\
    \texttt{ProdID2}         & Destination product of the transition. \\
    \texttt{mcatid1}         & Main-category identifier of the source product. \\
    \texttt{mcatid2}         & Main-category identifier of the destination product. \\
    \texttt{subcatid1}       & Sub-category identifier of the source product. \\
    \texttt{subcatid2}       & Sub-category identifier of the destination product. \\
    \texttt{Transition\_Type}& Nature of the navigation producing the transition \\
                             & (e.g.\ product-to-product, or originating from a search source). \\
    \bottomrule
  \end{tabular}
\end{table}

%--------------------------------------------------------------
\section{Product Description Dataset}
\label{sec:product-desc}

In addition to the behavioural data, a product description dataset provides
catalogue-level information about the products listed on IndiaMART. Unlike the
click stream log and the transition dataset, which describe what users did, this
dataset describes the products themselves and their sellers. It is the source of
product metadata used for content-based signals, feature construction, and
product representation learning in later chapters.

Each entry corresponds to a single catalogue listing and contains identifiers,
product information, seller details, category assignments, geographical
information, and product specifications. These attributes let a model reason
about products using their semantic properties in addition to how users interact
with them. The most informative of these attributes is the
\texttt{custom\_specs\_json} field, which stores product specifications as
structured JSON. These specifications capture characteristics of a product that
its categorical identifiers alone do not, and they can be turned into structured
features or into text for downstream models. Table~\ref{tab:product-schema}
lists the attributes of the dataset.

\begin{table}[ht]
  \centering
  \caption{Schema of the product description dataset.}
  \label{tab:product-schema}
  \begin{tabular}{ll}
    \toprule
    Attribute & Description \\
    \midrule
    \texttt{pc\_item\_display\_id}   & Unique identifier of the catalogue entry. \\
    \texttt{pc\_item\_name}          & Product name or title. \\
    \texttt{pc\_item\_img\_small}    & Reference to the product image for the entry. \\
    \texttt{pc\_item\_glusr\_usr\_id}& Identifier of the seller / catalogue owner. \\
    \texttt{city\_name}              & Seller's city. \\
    \texttt{state\_name}             & Seller's state. \\
    \texttt{fk\_glcat\_mcat\_id}     & Main product-category identifier. \\
    \texttt{glcat\_cat\_id}          & Product-category identifier. \\
    \texttt{glcat\_grp\_id}          & Product-group identifier. \\
    \texttt{prime\_pmcat\_id}        & Primary product category assigned to the entry. \\
    \texttt{custom\_specs\_json}     & Structured product specifications (attribute--value pairs) \\
                                     & stored as JSON. \\
    \bottomrule
  \end{tabular}
\end{table}

The product description dataset complements the behavioural data by describing
the intrinsic properties of catalogue items rather than user preferences. Where
the interaction logs reveal what users are drawn to, the catalogue metadata
records what each product is. Combining the two lets a recommender tell
genuinely similar products apart, build better item representations, and lean on
content when interaction data is sparse or a product is newly listed --- the
cold-start situations that recur throughout this work.

%--------------------------------------------------------------
\section{Data Preprocessing}
\label{sec:data-preprocessing}

Before any model is trained, the datasets are cleaned and put into a form the
learning algorithms can consume. This section covers the cleaning applied to the
transition dataset. The construction of content-based product vectors from the
description dataset is a preprocessing step too, but it is described in
Chapter~\ref{ch:embedding}, next to the embedding model that consumes those
vectors.

\subsection{Transition Dataset Cleaning}
\label{sec:transition-cleaning}
The transition dataset uses two special boundary markers to delimit browsing
sequences. A transition whose source is $\texttt{ProdID1} = -1$ marks the start
of a new user sequence, and a transition whose destination is
$\texttt{ProdID2} = -2$ marks the end of one. These values are artefacts
introduced when sequences are assembled; they do not correspond to any real
product on the platform.

Because the models are meant to learn relationships between actual products,
every row containing one of these boundary tokens is removed before training.
Keeping them would inject artificial transitions that do not reflect real
browsing behaviour and would distort the learned product representations. After
this filtering, the transition dataset contains only genuine
product-to-product movements, and it is this cleaned dataset that the models in
the following chapters are trained on.
